\section{The engineering, operations and reliability model}
\label{sec:model}

The model has two layers. A deterministic screen implements the protocol's closed-form
equations and reconciles them against the source workbook. A stochastic daily network
simulation over five measured years after a 365-day warm-up produces every service and cost
number reported. The screen audits arithmetic and bounds and never ranks architectures, because
its shortage-day column ignores safety stock and release queues. Apart from the four release
components marked tier 1 and 4 in Table~\ref{tab:release}, every input below is tier 5,
illustrative.

\subsection{Capacity, unit cost, response time, break-even}

Saleable annual capacity over $n$ sites with $q$ units per batch, $b$ nominal batches per
site-year, yield $y$ and uptime $u$, and capacity utilization against annual demand $D$, are
\begin{equation}
C = n \, q \, b \, y \, u ,
\qquad
U = D / C .
\label{eq:capacity}
\end{equation}
Yield and uptime enter the model only here, so a proportional change in either is
arithmetically the same lever as the same proportional change in $b$. Capital is annualized by
the capital recovery factor at discount rate $r$ over economic life $L$, with the limit $1/L$ at
$r=0$; fully loaded cost per delivered unit is the sum of the ten annual ledgers divided by
delivered non-expired units $Q$:
\begin{equation}
\mathrm{CRF}(r,L) = \frac{r(1+r)^{L}}{(1+r)^{L}-1} ,
\qquad
c = \frac{1}{Q}\bigl( K + F + V + M + O_{\mathrm{inv}} + T + W + I + R + G \bigr),
\label{eq:unitcost}
\end{equation}
where $K$ is annualized capital, $F$ fixed site operations including the quality unit and
labour, $V$ product and site launch, $M$ variable production, $O_{\mathrm{inv}}$ the opening
inventory purchase, $T$ quality testing, $W$ failure, scrap and expiry waste, $I$ inventory
carrying and distribution, $R$ resilience contracts and $G$ operating-system integration. The
denominator is delivered non-expired units, not demand: the workbook divided by demand, which
credits a plant for units it cannot make. Break-even committed volume, the resilience premium
over committed units, and incremental cost per shortage-day avoided against a reference are
\begin{equation}
Q_{\mathrm{be}} = \frac{F+R}{m} ,
\qquad
p = \frac{C_{s}-C_{0}}{Q_{\mathrm{committed}}} ,
\qquad
\mathrm{ICER} = \frac{C_{s}-C_{\mathrm{ref}}}{SD_{\mathrm{ref}}-SD_{s}} ,
\label{eq:economics}
\end{equation}
with $m$ the contribution margin per delivered unit, price net of variable production, testing,
distribution and expected waste, and $SD$ shortage days. The ratio is computed only where the
comparison is a genuine trade-off; a strategy that costs more and avoids no shortage days is
labelled dominated and receives none. Price, reimbursement, willingness to pay, provider
shortage cost and production cost are separate arguments in the economics module and are never
summed. CMS average sales price and Part B spending appear nowhere as a manufacturing cost.

Response time as the workbook sums it is serial with no overlap,
$T_{\mathrm{resp}} = \ell_{\mathrm{mat}} + t_{\mathrm{change}} + t_{\mathrm{cycle}} +
\phi\, t_{\mathrm{rel}} + t_{\mathrm{deliv}}$, over raw-material lead time, changeover,
production cycle, release time scaled by a strategy factor $\phi$, and delivery. The protocol
warns that these overlap only where overlap is operationally valid, so the stochastic engine
does not use that sum. It runs setup, fill, test and release as queues and reports time to first
response and time to stable recovery, the second being the first day followed by 30 consecutive
days at or above the daily shortage threshold of 0.95.

\subsection{The daily engine and the release path}

Each simulated day runs the protocol's nine-step order: expire lots first-expiry-first-out,
advance demand, advance disruption states, receive materials and released batches, serve
regional demand within a backorder window, start batches where the finished-goods position falls
below target and materials permit, replenish and trigger contract or transfer policies, accrue
the ten ledgers, then assert a mass balance in integer units,
\begin{equation}
I_{0} + P_{\mathrm{saleable}} = H + T_{\mathrm{transit}} + Q_{\mathrm{release}} + E + S ,
\label{eq:massbalance}
\end{equation}
opening stock plus saleable production equals on hand plus in transit plus in the release queue
plus expired plus served. The assertion caught a real leak on the first full-horizon run, a
shipment that expired in transit and was counted as served. Line occupancy is $365.25/b/u$ days
per batch, stretched by the reciprocal of the site's available capacity fraction that day and
floored at changeover plus production cycle, so uptime is a time stretch rather than a discount
on output.

Release time is not one number. Sterility incubation, environmental monitoring, assay and
endotoxin start together when the batch is filled, so the critical path is their maximum, and
quality-unit review is serial after them:
\begin{equation}
t_{\mathrm{rel}} = \max\{t_{\mathrm{ster}},\, t_{\mathrm{EM}},\, t_{\mathrm{assay}},\,
t_{\mathrm{endo}}\} + t_{\mathrm{QA}} .
\label{eq:release}
\end{equation}
At base values this is $\max\{14,7,5,2\}+2 = 16$ days, of which 14 are sterility incubation and
are not reducible by any analytical method the study models. Release scenario R1, which reduces
administrative time, can act only on the serial term, and R3, which could remove the assay
component, sits off the critical path unless sterility is shortened first. R3 cannot be
instantiated while gate G15 is UNCERTAIN. One defect is stated rather than hidden: the
release-time factor $\phi$ is read only by the deterministic screen and never by the stochastic
engine (MD-18), which is one of two reasons architecture S6 returns service numbers identical to
S5. The aggregate release inputs of 14 and 12 days in Table~\ref{tab:ops} and the 16-day
decomposition disagree; that is a known inconsistency between screen and engine, and the engine
uses Eq.~\ref{eq:release}.

\begin{table}[htbp]
\centering
\caption{Release-time components, from \texttt{config/global.yaml}. The first four are the only
inputs in the model above tier 5.}
\label{tab:release}
\footnotesize
\setlength{\tabcolsep}{5pt}
\begin{tabular}{@{}llccl@{}}
\toprule
Component (parameter id) & Path & Base (d) & Range (d) & Tier \\
\midrule
\texttt{sterility\_incubation\_days} & parallel & 14 & 14 to 18 & 1 \\
\texttt{environmental\_monitoring\_days} & parallel & 7 & 5 to 7 & 4 \\
\texttt{assay\_days} & parallel & 5 & 1 to 7 & 4 \\
\texttt{endotoxin\_days} & parallel & 2 & 0 to 3 & 4 \\
\texttt{qa\_review\_days} & serial & 2 & 0 to 4 & 5 \\
\texttt{release\_time\_cv} & dispersion & 0.20 & 0.10 to 0.40 & 5 \\
\texttt{investigation\_duration\_days} & added on deviation & 21 & 7 to 60 & 5 \\
\bottomrule
\end{tabular}
\end{table}

\subsection{Operational and cost parameters}

Table~\ref{tab:ops} holds the operational inputs the specification names, all tier 5. Node count
is not a single parameter: the network is four equal regions by protocol choice, with 3 and 8 as
sensitivity settings, and site count is a design variable the optimizer searches. Nodes may
exceed the region count and are placed round-robin, and a second node inside one region shares
that region's geography common-cause group, so redundancy inside a region buys less than
redundancy across regions.

Cost inputs are tier 5 on the same low, base, high convention. Per unit: materials 0.60 USD for
sodium bicarbonate and 0.55 for norepinephrine (0.25 to 1.50 and 0.20 to 1.50), conversion 0.60
for both (0.20 to 1.50), distribution 0.20 (0.05 to 0.80). Per batch: testing 30{,}000 USD
(10{,}000 to 100{,}000). Per site: fixed quality and labour 3.0 million USD per year (1.5 to
6.0), installed capital 40 million USD (15 to 100), validation and transfer 4 million USD once
(1 to 12). Globally: discount rate 0.10 (0.06 to 0.15) over a 10-year life (7 to 15), inventory
carrying 0.20 per year (0.12 to 0.30), and a shared operating layer at 0.30 million USD per
site-year (0.10 to 0.80) charged only outside release scenario R0. The last is an open defect
(MD-25), because every design-space architecture runs R0 and so declares a shared software group
carrying hazard with no cost line.

\begin{table}[htbp]
\centering
\caption{Operational parameters, base value with declared low and high range. All tier 5,
illustrative, from the two product configurations and \texttt{config/global.yaml}.}
\label{tab:ops}
\scriptsize
\setlength{\tabcolsep}{4pt}
\begin{tabular}{@{}lllll@{}}
\toprule
Parameter id & Units & Sodium bicarbonate & Norepinephrine & Tier \\
\midrule
\texttt{annual\_demand\_units} & million units/yr & 1.20 [0.60, 2.40] & 0.90 [0.50, 1.80] & 5 \\
\texttt{units\_per\_batch} & thousand units & 25 [10, 60] & 30 [10, 75] & 5 \\
\texttt{batches\_per\_site\_year\_nominal} & batches/site-yr & 45 [30, 70] & 45 [30, 70] & 5 \\
\texttt{yield\_fraction} & fraction & 0.96 [0.90, 0.99] & 0.96 [0.90, 0.99] & 5 \\
\texttt{uptime\_fraction} & fraction & 0.85 [0.70, 0.95] & 0.85 [0.70, 0.95] & 5 \\
\texttt{production\_cycle\_days} & days/batch & 3 [1, 7] & 3 [1, 7] & 5 \\
\texttt{changeover\_days} & days & 3 [1, 10] & 3 [1, 10] & 5 \\
\texttt{release\_time\_days} & days/batch & 14 [5, 28] & 12 [5, 28] & 5 \\
\texttt{material\_lead\_time\_days} & days & 90 [30, 240] & 90 [30, 240] & 5 \\
\texttt{shelf\_life\_months} & months & 24 [12, 36] & 24 [12, 30] & 5 \\
\texttt{expiry\_scrap\_fraction} & fraction & 0.020 [0.005, 0.080] & 0.020 [0.005, 0.080] & 5 \\
\texttt{delivery\_days} & days & 2 [1, 5] & 2 [1, 5] & 5 \\
\texttt{target\_safety\_stock\_days} & days & 60 [30, 180] & 60 [30, 180] & 5 \\
\texttt{yield\_sd} & fraction & 0.02 [0.01, 0.05] & 0.02 [0.01, 0.05] & 5 \\
\texttt{node\_commissioning\_days} & days & 730 [365, 1095] & 730 [365, 1095] & 5 \\
\texttt{capacity\_expansion\_days} & days & 270 [90, 730] & 270 [90, 730] & 5 \\
\bottomrule
\end{tabular}
\end{table}

\subsection{Reliability}
\label{sec:reliability}

Reliability is where the architectures separate, so failure processes are modelled explicitly
rather than assumed into a service level. Five mechanisms generate the exogenous world, all at
the tier-5 rates of Table~\ref{tab:reliability}. Daily regional demand is annual demand times
regional share divided by 365.25, times growth, seasonality, a routine multiplier and a shock
multiplier, rounded to integer units; the routine multiplier is a mean-one lognormal AR(1)
process at \texttt{demand\_cv} with lag-1 autocorrelation, so a bad week is correlated with the
week before it, and shocks arrive as a Poisson process hitting all four regions with a declared
probability and otherwise one, with lognormal magnitude and duration. Sites and suppliers each
draw Poisson failure arrivals with lognormal durations; degradation is proportional rather than
binary, so a line at capacity fraction $f$ takes $1/f$ times as long and a supplier at $f$ ships
at most $f$ of one full order per day. Since revision R011 pipeline cover is separated from the
policy cover of safety plus cycle stock, so shortening a lead time no longer shrinks the buffer,
though a residual remains and is not argued away, because under a forced outage a deferred order
lands when the supplier resumes whatever its nominal lead. A site that does not exist yet
contributes nothing but sunk capital, and since revision R008 starts its commissioning clock at
the first measured day rather than at simulation day zero; before that fix, three commissioning
parameters coincided with the warm-up, so any lead time of a year or less cost nothing inside
the measured window, and the architecture that had benefited most from the artifact was dual
sourcing.

Nominal site count is not independent redundancy, and the protocol's fair-comparison rules make
that a hard requirement. Every site and supplier declares its common-cause groups, and a site's
available capacity on a day is the minimum of its own and that of every group it belongs to.
Groups arrive at \texttt{common\_cause\_events\_per\_year} with lognormal durations, and impact
is a Beta draw, so an event degrades rather than switches off. Six kinds of group are
represented: API supplier, container and closure supplier, stopper and seal supplier, shared
software or model version, ownership and the quality unit that performs disposition, and
geography. The last two carry the study's scepticism about nominal redundancy. A shared quality
unit means one legal holder disposing batches from several sites, which is the situation for
most contract-manufacturing architectures whatever their site count; five design-space
architectures were configured claiming an independence their own regulatory-owner field denied,
and revision R007 added the shared quality group to all five. The same revision fixed a defect
that had made supplier-level groups structurally inert, since shipment was gated on capacity
strictly above zero while a group only ever sets capacity to $1-\text{impact}$, so a shared
key-starting-material tier could never delay a unit at any parameter value (MD-19). Fixing the
structure did not make the mechanism bite at the frozen hazard rate, which is reported as a null
rather than dropped.

\begin{table}[htbp]
\centering
\caption{Reliability parameters, base with declared range, all tier 5 illustrative, from
\texttt{config/global.yaml}. Durations are lognormal medians at dispersion
\texttt{disruption\_duration\_cv}.}
\label{tab:reliability}
\scriptsize
\setlength{\tabcolsep}{4pt}
\begin{tabular}{@{}llll@{}}
\toprule
Mechanism & Parameter id & Base & Range \\
\midrule
Routine demand & \texttt{demand\_cv} & 0.15 & 0.08 to 0.30 \\
 & \texttt{demand\_autocorrelation} & 0.30 & 0.00 to 0.70 \\
 & \texttt{annual\_demand\_growth} & 0.01 & $-0.02$ to 0.05 \\
Demand shocks & \texttt{demand\_shocks\_per\_year} & 0.40 & 0.10 to 1.00 \\
 & \texttt{demand\_shock\_multiplier} & 1.35 & 1.10 to 2.00 \\
 & \texttt{demand\_shock\_duration\_days} & 30 & 7 to 90 \\
 & \texttt{demand\_shock\_all\_regions\_probability} & 0.50 & 0.20 to 0.90 \\
Site failures & \texttt{site\_failures\_per\_site\_year} & 0.20 & 0.05 to 0.60 \\
 & \texttt{site\_failure\_duration\_days} & 45 & 7 to 180 \\
 & \texttt{site\_failure\_residual\_capacity} & 0.00 & 0.00 to 0.50 \\
Supplier disruptions & \texttt{supplier\_disruptions\_per\_supplier\_year} & 0.20 & 0.05 to 0.60 \\
 & \texttt{supplier\_disruption\_duration\_days} & 45 & 14 to 120 \\
Common cause & \texttt{common\_cause\_events\_per\_year} & 0.10 & 0.02 to 0.30 \\
 & \texttt{common\_cause\_duration\_days} & 60 & 14 to 240 \\
 & \texttt{common\_cause\_capacity\_impact} & 0.50 & 0.20 to 1.00 \\
 & \texttt{common\_cause\_capacity\_impact\_sd} & 0.15 & 0.05 to 0.30 \\
Quality & \texttt{deviation\_rate\_per\_batch} & 0.020 & 0.005 to 0.080 \\
 & \texttt{batch\_rejection\_rate} & 0.010 & 0.002 to 0.040 \\
Regulatory state & \texttt{shortage\_list\_resolution\_rate\_per\_year} & 0.30 & 0.10 to 1.00 \\
 & \texttt{shortage\_list\_relisting\_rate\_per\_year} & 0.10 & 0.02 to 0.40 \\
 & \texttt{bud\_503b\_days} & 90 & 45 to 180 \\
Logistics & \texttt{transport\_time\_cv} & 0.25 & 0.10 to 0.50 \\
 & \texttt{disruption\_duration\_cv} & 0.60 & 0.30 to 1.00 \\
\bottomrule
\end{tabular}
\end{table}

\subsection{The regulatory state machine}

The sixteen gates are binary feasibility constraints, not costs. A FAIL excludes an
architecture. UNCERTAIN permits simulation for information only and can never receive a
favourable decision class, and there is no penalty path an optimizer can pay through. As of the
current runs all sixteen are UNCERTAIN with no reviewer assigned, so eligibility is
NO\_CONCLUSION for every one of the twenty architectures. The 503B pathway is not a permanent
pathway and is not modelled as one. It is a time-varying legal state,
\begin{equation}
\mathrm{eligible}(t) = \mathrm{listed}(t) \wedge \mathrm{gates}_{\mathrm{static}}
 \wedge \mathrm{permission}(\mathrm{region}) ,
\label{eq:503b}
\end{equation}
where $\mathrm{listed}(t)$ is a two-state Markov chain on the FDA shortage list with hazards
\texttt{shortage\_list\_resolution\_rate\_per\_year} and
\texttt{shortage\_list\_relisting\_rate\_per\_year}, initialised from the frozen 2026-09-02
snapshot; $\mathrm{gates}_{\mathrm{static}}$ requires G08 and G11; and $\mathrm{permission}$
requires G12, with regions lacking state licensure removed from the service set. The predicate
is checked when a batch is compounded and again on each shipment and issue day, because the
statute requires the condition at compounding, at distribution and at dispensing. Section
\ref{sec:regulatory} states what that implies.

\subsection{The architectures, and what the model does not represent}

Table~\ref{tab:architectures} lists the twenty architectures. All are optimized to the same
target, minimizing expected annual cost subject to mean fill at or above 0.99 and a probability
of at least 0.90 that a run meets it, on common random numbers. Since revision R010 every frozen
comparator searches the same inventory space the design-space strategies declared for
themselves, because the earlier comparison measured search freedom as much as architecture:
added central capacity moved from infeasible to feasible on both products at roughly 60 percent
of its declared-space cost. The symmetric test, narrowing the design-space strategies to the
comparators' old space, has not been run and would be the honest completion of the check.

\begin{table}[htbp]
\centering
\caption{The twenty architectures. S0 to S7 are the frozen comparators, carrying the four
comparator classes the specification names plus centralized expansion and the 503B responder.
S8 to S19 are the design-space additions of revision R006.}
\label{tab:architectures}
\scriptsize
\setlength{\tabcolsep}{4pt}
\begin{tabular}{@{}lp{60mm}lp{48mm}@{}}
\toprule
Id & Architecture & Id & Architecture \\
\midrule
S0 & Status quo centralized supply & S10 & Split tenancy, two independent multi-product hosts \\
S1 & Additional safety stock & S11 & Bright-stock campaign offtake, national reserve \\
S2 & Dual source plus contracts & S12 & Contracted registered-capacity supply base \\
S3 & Reserved flexible contract capacity & S13 & Contracted base, rotated surge, reserve stock \\
S4 & Additional centralized capacity & S14 & Cooperative-financed registered second source \\
S5 & Distributed regional nodes, conventional release & S15 & Capacity-adequate presentation, no new plant \\
S6 & Distributed nodes plus validated release assurance & S16 & Reserve-triggered campaign network \\
S7 & 503B shortage response (time-varying legal state) & S17 & Rotating prequalified campaigns, three registered lines \\
S8 & Acquired registered line, split API and geography & S18 & Dual committed supply, contract-mandated sources \\
S9 & Dual finished-dose source, key-starting-material tier & S19 & Maintained alternate-source switching package \\
\bottomrule
\end{tabular}
\end{table}

The limits are recorded rather than described around. One product per network, so no campaign
scheduling or cost pooling across a portfolio; this is the largest unquantified distortion in
the cost comparison, because a node's fixed cost is borne by a single presentation at low
utilization (MD-15). Four equal regions, so three of the four allocation policies are
arithmetically the same function at the base case. No supplier allocation rule during scarcity.
Provider harm and clinical substitution are not monetized. And the screen's common-impact factor
is derived from topology for S8 to S19 but hand-set for S0 to S7, so no screen-level
common-impact comparison across the two arms is valid (MD-22).
